\section{教学过程设计：人机协同四阶段}

\subsection{阶段一：问题定义与AI辅助方案设计（8课时）}

\textbf{核心任务}：学生分组（4-5人/组），明确选址目标与约束条件，界定"生态友好型"内涵，识别关键利益相关方，初步确定评价因子及优先级。

\textbf{AI介入方式}：AI辅助生成结构化分析框架。

\textbf{Prompt示例}：
\begin{itemize}
\item "请为'生态友好型物流园区'GIS选址分析设计完整技术框架，包括：数据需求清单、分析步骤流程、评价指标体系、权重确定方法、结果验证方法。"
\item "在多因子加权评价中，常用的权重确定方法有哪些？各有什么优缺点？推荐适合本科教学的方法并说明步骤。"
\item "请提供国家层面生态保护区划分标准（自然保护区、生态红线、湿地保护范围、河湖蓝线等），哪些应设为'绝对排除区'？"
\end{itemize}

\textbf{教学要点}：引导学生批判性评估AI方案，结合地域实际调整指标权重。例如西南山区需提高地质灾害敏感性权重，东部平原则需更关注洪水风险。培养学生专业判断能力是本阶段关键。

\textbf{成果}：《物流园区选址分析方案》（含数据需求清单、技术路线、评价指标体系、权重方案）。

\subsection{阶段二：数据获取与AI辅助处理（8课时）}

\textbf{核心任务}：获取OSM道路网/DEM数据，完成投影转换、拓扑检查，将DEM转换为坡度图并按适宜性标准重分类为5级。

\textbf{AI介入方式}：AI生成数据获取与预处理Python代码，学生在Jupyter Notebook运行验证。

\textbf{Prompt示例}：
\begin{itemize}
\item "如何使用Python从OpenStreetMap获取某市道路网数据？请给出使用osmnx库的完整代码示例，包括安装、数据下载、范围指定、等级分类、本地保存。"
\item "请写Python代码实现：读取DEM $\rightarrow$ 计算坡度 $\rightarrow$ 按[0-5°)、[5-15°)、[15-25°)、[25-35°)、[>35°]重分类为5级 $\rightarrow$ 保存结果。使用GeoPandas或Rasterio库。"
\end{itemize}

\textbf{教学要点}：AI代码可能存在边界条件处理不当、投影转换错误等问题，学生需理解代码逻辑才能发现修正。强调数据质量评估与一致性处理的重要性。

\textbf{成果}：预处理后的矢量/栅格数据集 + 数据处理日志（含AI交互记录）。

\subsection{阶段三：空间建模与AI协同优化（12课时）}

\textbf{核心任务}：构建多因子加权评价模型，包括：各因子缓冲区/距离图生成、图层栅格化、AHP权重确定、加权叠加分析、敏感性分析。

\textbf{AI介入方式}：AI生成空间建模代码；提供"反事实分析"支持——快速生成对比方案，帮助理解不同参数设置对结果的影响。

\textbf{Prompt示例}：
\begin{itemize}
\item "请用ArcPy或GeoPandas写多因子加权叠加分析代码：输入5个重分类栅格（值域1-5），权重[0.25, 0.20, 0.20, 0.20, 0.15]，计算加权求和并分5级输出，需含完整注释。"
\item "假设生态保护区因子权重从0.25提高到0.35，预测选址结果如何变化？从地理学角度分析说明什么问题？"
\item "如何用成本距离（Cost Distance）分析替代直线距离，更准确地评估交通可达性？请给出Python代码示例。"
\end{itemize}

\textbf{反事实分析示例}：当生态保护区权重提高后，原30\%中适宜区变为低适宜区，选址重心向东南偏移约2km。这说明交通效率与生态保护存在矛盾，需决策者权衡。反事实分析使学生能在有限课时内进行更多情景模拟，深化对政策敏感性的理解。

\textbf{成果}：选址适宜性评价图（5级）、权重合理性论证、反事实分析报告。

\subsection{阶段四：成果表达与AI辅助科研素养培育（8课时）}

\textbf{核心任务}：完善选址分析报告，撰写规范的GIS专业研究报告，撰写AI辅助GIS实践反思报告。

\textbf{AI介入方式}：AI辅助生成专业报告框架与科研规范指导。

\textbf{Prompt示例}：
\begin{itemize}
\item "请提供一份规范的GIS选址分析报告结构框架，包括：摘要、关键词、引言、研究区概况、数据与方法、结果分析、讨论与结论、参考文献等部分的具体写作要点。"
\item "如何规范地撰写GIS分析报告中的地图制图说明？请提供地图要素（标题、图例、比例尺、指北针、坐标系统）规范表达的示例。"
\item "提供一个撰写《AI与地理实践》反思报告的结构框架，引导学生分析：AI在哪些环节提供有效帮助？产生哪些误导或困扰？人机协同边界在哪？对末来地理研究中AI应用的展望？"
\end{itemize}

\textbf{教学要点}：反思报告是课程关键环节，引导学生从效率、理解、局限、伦理四维度审视AI工具，形成客观全面的认知。重点培养学生严谨的科研态度和批判性思维。

\textbf{成果}：（1）选址分析报告（PDF，含规范地图、权重表、结果分析）；（2）《AI辅助GIS实践反思报告》（Word）。